\bigskip
\noindent Autres pistes de l'onglet Findings pour ce thème, dans des cotes qui n'ont pas de section ici :\par

\begin{encadre}
\begin{theoreme*}[Grothendieck, pages 42 à 45, excision pour deux cribles ouverts\piste{105-derivator-crible-excision}]\nonverifie
Soient $\mathbb D$ un dérivateur et $X$ une petite catégorie recouverte par deux cribles ouverts $U$ et $U'$. Si $U \cap U' \to U'$ est une $\mathbb D$-équivalence cohomologique, alors $U \to X$ en est une. C'est la moitié de l'axiome~W7 pour $W_{\mathbb D}$, tirée des seuls axiomes des dérivateurs.
\end{theoreme*}
\end{encadre}
{\small Les feuillets en esquissent une démonstration incomplète (page 45). L'identification de $\gamma(1,1)$ à $H^\bullet_{\mathbb D}(X, \xi)$ porte en marge \q{à vérifier (ou à mettre comme axiome sur les dérivateurs)}, et l'édition donne le pas qui manque. Mathlib n'a pas les dérivateurs. On se reportera aux notes de Grothendieck (\lot{105}{3}{42}).\par}

\begin{encadre}
\begin{theoreme*}[Grothendieck, pages 20 et 22, rectification d'un carré relatif\piste{107-relative-square-rectification}]
Dans une catégorie de modèles, soit un carré $pg = g'i$ où $i : X \to X'$ est une cofibration et $p : Y \to Y'$ une fibration, et soient des cylindres $X \otimes J$, $X' \otimes J$ reliés par $i \otimes J$, tels que la flèche de coin $X \otimes J \sqcup_X X' \to X' \otimes J$ soit une cofibration triviale. Soit $\varphi : X' \to Y$ une diagonale. Il existe une diagonale $\overline{\varphi}$, homotope à $\varphi$, avec $\overline{\varphi} i = g$ et $p\overline{\varphi} = g'$, si et seulement s'il existe des homotopies $k : \varphi i \to g$ et $h : p\varphi \to g'$ compatibles, c'est-à-dire $p \circ k = h \circ (i \otimes J)$. Si $Y'$ est l'objet final, la compatibilité est automatique.
\end{theoreme*}
\end{encadre}
{\small \emph{Démonstration.} Suffisance : les données $\varphi$ sur $X' \otimes \{0\}$ et $k$ sur $X \otimes J$ se recollent en une flèche $X \otimes J \sqcup_X X' \to Y$ ; la compatibilité dit exactement que, composée avec $p$, elle est la restriction de $h$. Le carré ainsi formé a une cofibration triviale à gauche et la fibration $p$ à droite : il a un relèvement $L : X' \otimes J \to Y$, homotopie de $\varphi$ à $\overline{\varphi} = L \circ i_1$. Alors $\overline{\varphi} i = L \circ (i \otimes J) \circ i_1 = k \circ i_1 = g$ et $p\overline{\varphi} = h \circ i_1 = g'$. Nécessité : une homotopie $K : \varphi \to \overline{\varphi}$ donne $k = K \circ (i \otimes J)$ et $h = p \circ K$, qui sont compatibles par associativité. Si $Y'$ est final, deux flèches vers $Y'$ sont égales. Lean (\href{\leanblob/Folder107.lean}{\texttt{Folder107.lean}}) démontre la suffisance (\texttt{rectification}) sous la seule propriété de relèvement de $i \otimes J$ contre $p$, la nécessité (\texttt{compatible\_of\_rectification}), le critère (\texttt{rectification\_iff}, et \texttt{rectification\_modelCategory} dans une \texttt{ModelCategory} de Mathlib, où la propriété vient de l'axiome de relèvement) et le cas final (\texttt{rectification\_of\_isTerminal}).\par}
{\small La compatibilité est prise au sens strict de la page~22. Ne sont pas démontrés : la forme « à homotopie relative aux extrémités près » de la marge de la page~20, lecture de l'édition, qui demande des cylindres de cylindres ; l'existence, dans une catégorie de modèles quelconque, de cylindres relatifs dont la flèche de coin est une cofibration triviale (axiome SM7 dans le cas simplicial), qui est ici une hypothèse ; et le contre-exemple de l'édition (fibration des chemins sur $S^1$) montrant que la condition peut manquer. On se reportera aux notes de Grothendieck (\lot{107}{1}{20}).\par}

\begin{encadre}
\begin{theoreme*}[Grothendieck, page 2, $\Delta^\circ$ et les catégories test\piste{108-delta-op-not-test}]\nonverifie
La catégorie $\Delta^\circ$ n'est pas une catégorie test.
\end{theoreme*}
\end{encadre}
{\small La question~6 de la liste reçoit pour réponse un \q{non} souligné, sans argument ; la page ne définit ni modélisateur ni catégorie test, et l'édition ne reconstruit pas d'argument. Mathlib n'a pas les catégories test. On se reportera aux notes de Grothendieck (\lot{108}{1}{2}).\par}

\begin{encadre}
\begin{theoreme*}[D'après les pages 18 à 20, forme complétée par l'édition, sommes amalgamées le long d'opfibrations ; voir aussi Cisinski~\cite{cisinski2006}\piste{118-pushout-along-cofibrations-in-cat}]\nonverifie
Soient $\mathcal W$ un localisateur fondamental et, dans $\mathbf{Cat}$, un carré cocartésien et cartésien formé de $i : Y \to X$, $g : Y \to Y'$, $f : X \to X'$ et $i' : Y' \to X'$, où $i'$ et $f$ sont des opfibrations de Grothendieck et où, au-dessus de chaque objet de $X'$, l'une des fibres de $f$ ou de $i'$ est équivalente au point. Ce carré est homotopiquement cocartésien : si $i \in \mathcal W$, alors $i' \in \mathcal W$, et si $g \in \mathcal W$, alors $f \in \mathcal W$.
\end{theoreme*}
\end{encadre}
{\small Les feuillets en esquissent une démonstration incomplète (pages 18 à 20). L'argument de cocartésianité s'interrompt, et l'hypothèse cartésienne est de la lecture. Cisinski conditionne ses énoncés voisins à des immersions, non à des opfibrations. Mathlib n'a pas les localisateurs fondamentaux. On se reportera aux notes de Grothendieck (\lot{118}{1}{18}).\par}

\begin{encadre}
\begin{theoreme*}[Grothendieck, pages 161 à 179 et 189 à 199, schéma en groupes des automorphismes\piste{136-automorphism-group-scheme}]\nonverifie
Les homomorphismes d'anneaux simpliciaux entre algèbres de de Rham à puissances divisées $\mathcal A(S, S^+, t)_*$, sans graduation, filtration ni puissances divisées imposées, sont classés par des données $(\varphi_0, A'_1, (B'_i))$ soumises à trois conditions, la troisième étant $(2B'_i)J' = 0$. Le foncteur des automorphismes de l'algèbre complétée sur $k\{T\}$ est donc un schéma en groupes affine sur $\mathbf Z$, produit semi-direct de $\mathbf G_m$ par un groupe unipotent dont les facteurs sont tués par~$2$. Quand $2$ est inversible, tout automorphisme respecte la graduation extérieure et opère sur la cohomologie par $\mathbf G_m$ seul.
\end{theoreme*}
\end{encadre}
{\small Les feuillets en esquissent une démonstration incomplète (pages 173 à 179) : la condition~3) est affirmée sans raison, et la raison qu'en donne la lecture est la sienne. Mathlib n'a pas les algèbres simpliciales à puissances divisées. On se reportera aux notes de Grothendieck (\lot{136}{9}{161}).\par}

\begin{encadre}
\begin{theoreme*}[Grothendieck, pages 5 à 10 et 226 à 232, ombres de la cohomologie entière ; voir aussi Kageyama 2022\piste{136-pd-de-rham-shadows}]\nonverifie
Soit $C^{p,q} = \Gamma^p\Phi \otimes \Lambda^q\Psi$ la $k\{T\}$-algèbre simpliciale des formes à puissances divisées sur le simplexe $\sum X_i = T$. Pour tout ensemble simplicial $X$, la cohomologie de ses sections sur $X$ est $H^q(X, k)$ en tout degré complémentaire $p \geq 0$ ; graduée par le degré total, $H^{\bullet,q}$ est la troncature $\tau_{\geq q}(H^q(X, k) \otimes_k k\{T\})$. La cohomologie entière n'est retrouvée que sous la forme de ces ombres.
\end{theoreme*}
\end{encadre}
{\small La page~8 affirme sans démonstration l'exactitude de chaque ligne ; la lecture la démontre à partir du complexe de Koszul à puissances divisées, et les pages 226 à 232 refont le calcul. Kageyama construit la même algèbre sans en calculer la cohomologie. Mathlib n'a pas la cohomologie des ensembles simpliciaux à coefficients dans ces algèbres. On se reportera aux notes de Grothendieck (\lot{136}{1}{5}).\par}

\begin{encadre}
\begin{theoreme*}[Grothendieck, pages 10 à 14 et 17, hyperrecouvrements indexés par une catégorie quelconque\piste{151-hypercovering-any-index-category}]\nonverifie
Un hyperrecouvrement d'un topos $X$ indexé par une petite catégorie quelconque $A$, et non par un objet simplicial, est caractérisé par un seul axiome : la flèche de comparaison $\varphi_!(\varprojlim_I F_i) \to \varprojlim_I \varphi_!(F_i)$ est un épimorphisme pour tout diagramme fini de $\hat A$. La version où elle est un isomorphisme équivaut à l'exactitude à gauche de $\varphi_!$, c'est-à-dire à un morphisme de topos $X \to \hat A$.
\end{theoreme*}
\end{encadre}
{\small Les quatre axiomes des pages 10 à 13 se ramènent à celui-ci par $F = \varinjlim_{A/F} \alpha$ (pages 13, 14 et 17). Mathlib a les préfaisceaux et les extensions de Kan, mais pas les hyperrecouvrements d'un topos. On se reportera aux notes de Grothendieck (\lot{151}{1}{10}).\par}

\begin{encadre}
\begin{definition*}[Pages 59, 60 et 74, forme corrigée par l'édition, tubes mixtes\piste{151-mixed-tubes}]
Pour $i \leq j \leq k$, le tube mixte $V^j_{i,k}$ s'obtient en retirant du tube de $X_i$ dans $X_k$ la réunion $R_{i,k} \cup S_{j,k}$ de deux parties complémentaires du bord de $X_k$. Ces tubes interpolent de façon monotone entre le tube épointé $V^*_{i,k}$, pour $j = i$, et le tube plein $V_{i,k}$, pour $j = k$.
\end{definition*}
\end{encadre}
{\small La page écrit $R_{j,k} \cup S_{j,k}$, qui est le bord de $X_k$ pour tout $j$ : la famille serait constante, et l'indice $i$ est de l'édition. La page~74 commence une généralisation indexée par des cribles et s'interrompt. Mathlib n'a pas les espaces stratifiés ni leurs tubes. On se reportera aux notes de Grothendieck (\lot{151}{3}{59}).\par}

\begin{encadre}
\begin{exemple*}[Grothendieck, pages 61 à 63, $\Pi_1$ trop grossier pour la sphère\piste{151-pi1-too-coarse-for-the-sphere}]\nonverifie
Sur $\mathbf P^1_{\mathbf C} \simeq S^2$, stratifié par un point, un plan et un cercle, la somme amalgamée du diagramme des tubes au niveau du groupoïde fondamental est homotopiquement triviale, alors que la sphère ne l'est pas. En montant d'un niveau, aux gerbes, on retrouve $H^2(S^2, A) \simeq A$.
\end{exemple*}
\end{encadre}
{\small Mathlib n'a pas la théorie de l'homotopie des espaces stratifiés ni les gerbes. On se reportera aux notes de Grothendieck (\lot{151}{4}{61}).\par}

\begin{encadre}
\begin{theoreme*}[Grothendieck, pages 62 à 68, recollement par les paires de couverture\piste{151-tube-gluing-covering-pairs}]\nonverifie
Un espace stratifié est la limite inductive d'un diagramme indexé par les seules paires $(i, j)$ de l'ensemble ordonné des strates où $j$ couvre $i$, avec trois flèches par paire, $X_i^* \to V_{i,j} \leftarrow V^*_{i,j} \to X_j^*$ ; aucune comparaison entre strates non adjacentes n'est nécessaire.
\end{theoreme*}
\end{encadre}
{\small Les feuillets en esquissent une démonstration incomplète (pages 62 et 63) : ils établissent, sous trois hypothèses, que ces trois inclusions sont les seules entre les pièces élémentaires. Le théorème de recollement de la page~68 est sans démonstration, et le mot qui y nomme l'opération est illisible. Mathlib n'a pas les espaces stratifiés. On se reportera aux notes de Grothendieck (\lot{151}{4}{62}).\par}
